% macros.tex
% 课程语义命令、常用公式命令与复习资料专用环境。

% ---------- 文档信息 ----------
\newcommand{\USTC}{中国科学技术大学}
\newcommand{\CourseName}{热学B}
\newcommand{\DocTitle}{\USTC《\CourseName》复习资料}
\newcommand{\DocVersion}{2026.1}

% 作者信息
\newcommand{\DocAuthor}{寒初晓}
\newcommand{\DocAuthorInfo}{中国科学技术大学 \quad 近代物理系}

% ---------- 封面 ----------
\newcommand{\makecoursecover}{%
  \begin{titlepage}
    \thispagestyle{empty}

    \vspace*{1.9cm}

    \begin{flushleft}
      {\color{USTCBlue}\rule{0.20\textwidth}{2.2pt}\par}

      \vspace{1.35cm}

      {\zihao{-4}\color{Gray}\textsf{THERMAL PHYSICS B}\par}

      \vspace{0.55cm}

      {\zihao{0}\bfseries 中国科学技术大学\par}

      \vspace{0.35cm}

      {\zihao{1}\bfseries 《热学B》复习资料\par}
    \end{flushleft}

    \begin{flushleft}
      \renewcommand{\arraystretch}{1.55}
      \begin{tabular}{@{}p{0.20\textwidth}p{0.70\textwidth}@{}}
        \textcolor{Gray}{版本} &
        \DocVersion \\
      \end{tabular}
    \end{flushleft}

    \vfill

    \begin{flushleft}
      {\color{USTCBlue}\rule{\textwidth}{0.7pt}\par}

      \vspace{0.45cm}

      {\zihao{5}\color{Gray}
      \DocAuthor
      \par}

      \vspace{0.15cm}

      {\zihao{5}\color{Gray}
      \DocAuthorInfo
      \par}

      \vspace{0.15cm}

      {\zihao{5}\color{Gray}
      \today
      \par}
    \end{flushleft}

  \end{titlepage}
}

% ---------- 微分、偏导与过程量 ----------
\newcommand{\dd}{\mathop{}\!\mathrm{d}}
\newcommand{\dbar}{\mathop{}\!\bar{\mathrm{d}}}

\newcommand{\pd}[2]{\frac{\partial #1}{\partial #2}}
\newcommand{\pdc}[3]{\left(\frac{\partial #1}{\partial #2}\right)_{#3}}

\newcommand{\od}[2]{\frac{\dd #1}{\dd #2}}
\newcommand{\evalat}[2]{\left.#1\right|_{#2}}

% ---------- 常用热学量 ----------
\newcommand{\kB}{k_{\mathrm B}}
\newcommand{\NA}{N_{\mathrm A}}

\newcommand{\Cv}{C_V}
\newcommand{\Cp}{C_p}
\newcommand{\Cvm}{C_{V,m}}
\newcommand{\Cpm}{C_{p,m}}

\newcommand{\gam}{\gamma}

\newcommand{\avg}[1]{\left\langle #1 \right\rangle}

\newcommand{\vrms}{v_{\mathrm{rms}}}
\newcommand{\vmp}{v_{\mathrm p}}
\newcommand{\vmean}{\bar v}

\newcommand{\eps}{\varepsilon}

% ---------- 复习资料语义环境 ----------
\newenvironment{ReviewList}
  {\begin{enumerate}[label=\textbf{\arabic*.},leftmargin=2.5em]}
  {\end{enumerate}}

\newenvironment{ConditionList}
  {\begin{itemize}[label=\textcolor{USTCBlue}{\ensuremath{\blacktriangleright}},leftmargin=2.2em]}
  {\end{itemize}}

\newenvironment{PitfallList}
  {\begin{itemize}[label=\textcolor{WarnRed}{\ensuremath{\times}},leftmargin=2.2em]}
  {\end{itemize}}

\newcommand{\concept}[1]{\textbf{\color{USTCBlue}#1}}
\newcommand{\keyword}[1]{\textbf{#1}}
\newcommand{\condition}[1]{\textit{适用条件：#1}}
\newcommand{\exam}[1]{\textcolor{WarnRed}{\textbf{考点：}#1}}
\newcommand{\sourcepdf}[1]{\par\smallskip\noindent\textcolor{Gray}{\small 来源参考：#1}}

\newcommand{\notA}{\textbf{注意：本资料按《热学B》范围整理，不按《热学A》额外拔高。}}

% ---------- 课程核心公式快捷命令 ----------

% 热力学基础
\newcommand{\IdealGas}{pV=\nu RT=Nk_{\mathrm B}T}
\newcommand{\VanDerWaals}{\left(p+\dfrac{\nu^2 a}{V^2}\right)(V-\nu b)=\nu RT}

% 第一章、第二章常用
\newcommand{\FirstLaw}{\dd U=\dbar Q+\dbar W}
\newcommand{\QuasiWork}{\dbar W=-p\dd V}
\newcommand{\HeatCapacityDef}{C_L=\left(\dfrac{\dbar Q}{\dd T}\right)_L}
\newcommand{\EnthalpyDef}{H=U+pV}

% 理想气体过程
\newcommand{\AdiabaticIdeal}{
  pV^{\gamma}=\mathrm{const},\quad
  TV^{\gamma-1}=\mathrm{const},\quad
  p^{1-\gamma}T^{\gamma}=\mathrm{const}
}

% 第二定律与熵
\newcommand{\EntropyDef}{\dd S=\dfrac{\dbar Q_{\mathrm{rev}}}{T}}
\newcommand{\ClausiusIneq}{\oint\dfrac{\dbar Q}{T}\le 0}
\newcommand{\BoltzmannEntropy}{S=\kB\ln W}

% 麦克斯韦--玻尔兹曼分布
\newcommand{\PressureKinetic}{p=\dfrac{1}{3}nm\avg{v^2}=\dfrac{2}{3}n\avg{\varepsilon_k}}
\newcommand{\MeanKinetic}{\avg{\varepsilon_k}=\dfrac{3}{2}\kB T}

% 输运过程
\newcommand{\FourierLaw}{j_Q=-\kappa\od{T}{z}}

% 相变与液体表面性质
\newcommand{\Clapeyron}{\od{p}{T}=\dfrac{l}{T(v_2-v_1)}}
\newcommand{\CapillaryRise}{h=\dfrac{2\sigma\cos\theta}{\rho g r}}

% ---------- 公式卡片快捷结构 ----------
\newcommand{\FormulaCard}[4]{%
  \begin{FormulaBox}[#1]
  \begin{align*}
    #2
  \end{align*}
  \condition{#3}
  \par\smallskip
  \textbf{复习提示：}#4
  \end{FormulaBox}
}

\newcommand{\Pitfall}[2]{%
  \begin{WarnBox}[#1]
  #2
  \end{WarnBox}
}

\newcommand{\Method}[2]{%
  \begin{MethodBox}[#1]
  #2
  \end{MethodBox}
}

% ---------- 占位命令：方便后续批量写章节 ----------
\newcommand{\chaptergoal}[1]{%
  \begin{KeyBox}[本章目标]
  #1
  \end{KeyBox}
}

\newcommand{\exammap}[1]{%
  \begin{MethodBox}[考点地图]
  #1
  \end{MethodBox}
}

\newcommand{\formulatable}[1]{%
  \begin{FormulaBox}[公式速查]
  #1
  \end{FormulaBox}
}

% ---------- 第二章：热力学第一定律常用命令 ----------

% 功的两种符号习惯：本资料默认 W 为外界对系统做功
\newcommand{\WorkOnSystem}{\dbar W=-p\dd V}
\newcommand{\WorkBySystem}{\dbar W_{\mathrm{by}}=p\dd V}

% 第一定律与焓
\newcommand{\FirstLawFull}{\dd U=\dbar Q+\dbar W=\dbar Q-p\dd V}
\newcommand{\Mayer}{C_{p,m}-C_{V,m}=R}

% 理想气体内能、焓
\newcommand{\IdealDU}{\dd U=\nu C_{V,m}\dd T}
\newcommand{\IdealDH}{\dd H=\nu C_{p,m}\dd T}

% 理想气体等温、绝热、多方过程
\newcommand{\IsoThermalWorkOn}{W=-\nu RT\ln\frac{V_2}{V_1}}
\newcommand{\AdiabaticWorkOn}{W=\Delta U=\nu C_{V,m}(T_2-T_1)}
\newcommand{\PolyProcess}{pV^n=\mathrm{const}}
\newcommand{\PolyHeatCapacity}{C_{n,m}=C_{V,m}\frac{\gamma-n}{1-n}}

% 循环、热机、制冷机
\newcommand{\EngineEff}{\eta=\frac{W'}{Q_1}=1-\frac{Q_2}{Q_1}}
\newcommand{\CarnotEff}{\eta_{\mathrm C}=1-\frac{T_2}{T_1}}
\newcommand{\COP}{\varepsilon=\frac{Q_2}{W}=\frac{Q_2}{Q_1-Q_2}}
\newcommand{\CarnotCOP}{\varepsilon_{\mathrm C}=\frac{T_2}{T_1-T_2}}

% 焦耳--汤姆孙系数
\newcommand{\JTcoef}{\alpha_i=\left(\frac{\partial T}{\partial p}\right)_H}

% ---------- 第三章：热力学第二定律与熵常用命令 ----------

% 第二定律、克劳修斯关系与熵
\newcommand{\ClausiusEq}{\oint_{\mathrm{rev}}\frac{\dbar Q}{T}=0}
\newcommand{\ClausiusGeneral}{\oint\frac{\dbar Q}{T}\le 0}
\newcommand{\EntropyIntegral}{\Delta S=S_b-S_a=\int_a^b\frac{\dbar Q_{\mathrm{rev}}}{T}}
\newcommand{\EntropyIneq}{\dd S\ge \frac{\dbar Q}{T}}
\newcommand{\IsolatedEntropy}{\Delta S_{\mathrm{isolated}}\ge 0}

% 理想气体熵变
\newcommand{\IdealEntropyTV}{\Delta S=\nu C_{V,m}\ln\frac{T_2}{T_1}+\nu R\ln\frac{V_2}{V_1}}
\newcommand{\IdealEntropyTp}{\Delta S=\nu C_{p,m}\ln\frac{T_2}{T_1}-\nu R\ln\frac{p_2}{p_1}}
\newcommand{\IdealEntropypV}{\Delta S=\nu C_{V,m}\ln\frac{p_2}{p_1}+\nu C_{p,m}\ln\frac{V_2}{V_1}}

% 卡诺热机、卡诺制冷机
\newcommand{\CarnotHeatRatio}{\frac{Q_2}{Q_1}=\frac{T_2}{T_1}}
\newcommand{\TSarea}{W'=\oint T\dd S}

% ---------- 第四章：麦克斯韦--玻尔兹曼分布律常用命令 ----------

% 麦克斯韦速度与速率分布
\newcommand{\MaxwellComponent}{
  f(v_x)=\left(\frac{m}{2\pi k_{\mathrm B}T}\right)^{1/2}
  \exp\left(-\frac{mv_x^2}{2k_{\mathrm B}T}\right)
}

\newcommand{\MaxwellVelocity}{
  f(v_x,v_y,v_z)=
  \left(\frac{m}{2\pi k_{\mathrm B}T}\right)^{3/2}
  \exp\left[-\frac{m(v_x^2+v_y^2+v_z^2)}{2k_{\mathrm B}T}\right]
}

\newcommand{\MaxwellSpeed}{
  f(v)=4\pi
  \left(\frac{m}{2\pi k_{\mathrm B}T}\right)^{3/2}
  v^2
  \exp\left(-\frac{mv^2}{2k_{\mathrm B}T}\right)
}

\newcommand{\MostProbSpeed}{
  v_{\mathrm p}=\sqrt{\frac{2k_{\mathrm B}T}{m}}
  =\sqrt{\frac{2RT}{\mu}}
}

\newcommand{\MeanSpeed}{
  \bar v=\sqrt{\frac{8k_{\mathrm B}T}{\pi m}}
  =\sqrt{\frac{8RT}{\pi\mu}}
}

\newcommand{\RmsSpeed}{
  v_{\mathrm{rms}}=\sqrt{\overline{v^2}}
  =\sqrt{\frac{3k_{\mathrm B}T}{m}}
  =\sqrt{\frac{3RT}{\mu}}
}

% 碰壁数与泻流
\newcommand{\CollisionRate}{
  G=\frac{1}{4}n\bar v
}

\newcommand{\EffusionRate}{
  \frac{\Delta N}{\Delta t}=\frac{1}{4}n\bar v A
}

% 玻尔兹曼分布
\newcommand{\BoltzmannDist}{
  n(\bm r)=n_0\exp\left[-\frac{\varepsilon_p(\bm r)}{k_{\mathrm B}T}\right]
}

\newcommand{\BarometricFormula}{
  n(z)=n_0\exp\left(-\frac{mgz}{k_{\mathrm B}T}\right),
  \qquad
  p(z)=p_0\exp\left(-\frac{\mu gz}{RT}\right)
}

% 能量均分
\newcommand{\Equipartition}{
  \overline{\varepsilon_i}=\frac{1}{2}k_{\mathrm B}T
}

\newcommand{\MolecularEnergy}{
  \bar\varepsilon=\frac{i}{2}k_{\mathrm B}T
}

\newcommand{\IdealCvFromFreedom}{
  C_{V,m}=\frac{i}{2}R
}

% ---------- 第五章：气体中的输运过程常用命令 ----------

% 宏观输运规律
\newcommand{\FourierTransport}{
  \frac{\Delta Q}{\Delta t}
  =
  -\kappa\frac{\dd T}{\dd z}\Delta S
}

\newcommand{\NewtonViscosity}{
  f
  =
  -\eta\frac{\dd u}{\dd z}\Delta S
}

\newcommand{\FickLaw}{
  \frac{\Delta N}{\Delta t}
  =
  -D\frac{\dd n}{\dd z}\Delta S
}

% 输运系数的微观表达
\newcommand{\ThermalConductivity}{
  \kappa=\frac13\lambda\rho c_V\bar v
}

\newcommand{\ViscosityCoeff}{
  \eta=\frac13\lambda\rho\bar v
}

\newcommand{\DiffusionCoeff}{
  D=\frac13\lambda\bar v
}

% 平均自由程与碰撞频率
\newcommand{\MeanFreePath}{
  \lambda=\frac{1}{\sqrt{2}\pi d^2 n}
}

\newcommand{\CollisionFreq}{
  Z=\frac{\bar v}{\lambda}
  =
  \sqrt{2}\pi d^2n\bar v
}

% 输运系数与压强、温度的近似关系
\newcommand{\KappaPressureIndep}{
  \kappa\sim \frac13\lambda\rho c_V\bar v
}

\newcommand{\EtaPressureIndep}{
  \eta\sim \frac13\lambda\rho\bar v
}

\newcommand{\DPressureRelation}{
  D\sim \frac13\lambda\bar v
}

% ---------- 第六章：固、液体性质简介与相变常用命令 ----------

% 固体热容量
\newcommand{\DulongPetit}{
  C_{V,m}\approx 3R
}

% 表面张力
\newcommand{\SurfaceForce}{
  f=\sigma l
}

\newcommand{\SurfaceWork}{
  \Delta A=\sigma\Delta S
}

% 弯曲液面附加压强
\newcommand{\LaplaceSphere}{
  \Delta p=\frac{2\sigma}{R}
}

\newcommand{\LaplaceCylinder}{
  \Delta p=\frac{\sigma}{R}
}

% 毛细现象
\newcommand{\CapillaryFormula}{
  h=\frac{2\sigma\cos\theta}{\rho gr}
}

% 克拉珀龙方程
\newcommand{\ClapeyronEq}{
  \frac{\dd p}{\dd T}=\frac{l}{T(v_2-v_1)}
}

% 克劳修斯--克拉珀龙近似
\newcommand{\ClausiusClapeyron}{
  \frac{\dd p}{\dd T}=\frac{lp}{RT^2}
}

\newcommand{\ClausiusClapeyronInt}{
  \ln p=-\frac{l}{RT}+C
}

% ---------- 总结部分常用命令 ----------

\newcommand{\finalnote}[1]{%
  \begin{KeyBox}[最后提醒]
  #1
  \end{KeyBox}
}